\textbf{Phần 1. Trắc nghiệm lựa chọn}
\setcounterpageref{ex}{0}

\begin{ex} %[5D7V40]
Theo kế hoạch, một xã dự định trồng $20\text{ ha}$ rừng. Thực tế xã đã trồng được $25\text{ ha}$ rừng. Xã đó đã vượt mức kế hoạch đề ra là bao nhiêu phần trăm?
\choice
{$5\%$}
{$20\%$}
{\True $25\%$}
{$125\%$}
\loigiai{
Tỉ số phần trăm diện tích rừng trồng được so với kế hoạch là:
$25 : 20 = 1{,}25 = 125\%$.\\
Số phần trăm vượt mức kế hoạch là:
$125\% - 100\% = 25\%$.\\
\textbf{Đáp án của câu hỏi là C}
}
\end{ex}

\begin{ex} %[5D7V40]
Giá của một chiếc quạt điện tăng từ $400\,000$ đồng lên $460\,000$ đồng. Giá của chiếc quạt điện đó đã tăng thêm bao nhiêu phần trăm?
\choice
{$12\%$}
{\True $15\%$}
{$18\%$}
{$60\%$}
\loigiai{
Số tiền tăng thêm là:
$460\,000 - 400\,000 = 60\,000\text{ (đồng)}$.\\
Tỉ số phần trăm giá tăng thêm so với giá ban đầu là:
$60\,000 : 400\,000 = 0{,}15 = 15\%$.\\
\textbf{Đáp án của câu hỏi là B}
}
\end{ex}

\begin{ex} %[5D7V40]
Một vườn cây có $500$ cây ăn quả, trong đó có $175$ cây cam. Tỉ số phần trăm của số cây còn lại so với tổng số cây trong vườn là
\choice
{$35\%$}
{\True $65\%$}
{$55\%$}
{$70\%$}
\loigiai{
Số cây còn lại trong vườn là:
$500 - 175 = 325\text{ (cây)}$.\\
Tỉ số phần trăm của số cây còn lại là:
$325 : 500 = 0{,}65 = 65\%$.\\
\textbf{Đáp án của câu hỏi là B}
}
\end{ex}

\begin{ex} %[5D7V40]
Một người gửi tiết kiệm $100\,000\,000$ đồng tại ngân hàng. Sau một năm, người đó nhận về cả gốc và lãi là $106\,500\,000$ đồng. Lãi suất tiết kiệm theo năm của ngân hàng là
\choice
{$6\%$}
{\True $6{,}5\%$}
{$7\%$}
{$7{,}5\%$}
\loigiai{
Số tiền lãi sau một năm là:
$106\,500\,000 - 100\,000\,000 = 6\,500\,000\text{ (đồng)}$.\\
Lãi suất tiết kiệm một năm là:
$6\,500\,000 : 100\,000\,000 = 0{,}065 = 6{,}5\%$.\\
\textbf{Đáp án của câu hỏi là B}
}
\end{ex}

\begin{ex} %[5D7V40]
Một mảnh vườn hình chữ nhật có chiều dài $20\text{ m}$ và chiều rộng $16\text{ m}$. Tỉ số phần trăm của chiều rộng so với chiều dài mảnh vườn là
\choice
{$75\%$}
{\True $80\%$}
{$85\%$}
{$125\%$}
\loigiai{
Tỉ số phần trăm của chiều rộng so với chiều dài là:
$16 : 20 = 0{,}8 = 80\%$.\\
\textbf{Đáp án của câu hỏi là B}
}
\end{ex}

\begin{ex} %[5D7V40]
Một cửa hàng nhập một chiếc máy tính xách tay với giá $15\,000\,000$ đồng và bán ra với giá $17\,250\,000$ đồng. Cửa hàng đó đã lãi bao nhiêu phần trăm so với giá vốn?
\choice
{$12\%$}
{\True $15\%$}
{$18\%$}
{$20\%$}
\loigiai{
Số tiền lãi là:
$17\,250\,000 - 15\,000\,000 = 2\,250\,000\text{ (đồng)}$.\\
Tỉ số phần trăm tiền lãi so với giá vốn là:
$2\,250\,000 : 15\,000\,000 = 0{,}15 = 15\%$.\\
\textbf{Đáp án của câu hỏi là B}
}
\end{ex}

\begin{ex} %[5D7V40]
Khối 5 của một trường tiểu học có $250$ học sinh, trong kì thi học kì 1 có $195$ bạn đạt điểm Giỏi. Tỉ số phần trăm số học sinh đạt điểm Giỏi là
\choice
{$75\%$}
{$76\%$}
{\True $78\%$}
{$80\%$}
\loigiai{
Tỉ số phần trăm học sinh đạt điểm Giỏi là:
$195 : 250 = 0{,}78 = 78\%$.\\
\textbf{Đáp án của câu hỏi là C}
}
\end{ex}

\begin{ex} %[5D7V40]
Nước biển chứa $3{,}5\%$ muối theo khối lượng. Để thu được $70\text{ kg}$ muối tinh, người ta cần làm bay hơi hoàn toàn bao nhiêu ki-lô-gam nước biển?
\choice
{$1\,500\text{ kg}$}
{\True $2\,000\text{ kg}$}
{$2\,450\text{ kg}$}
{$3\,500\text{ kg}$}
\loigiai{
Khối lượng nước biển cần làm bay hơi là:
$70 : 3{,}5\% = 70 : 0{,}035 = 2\,000\text{ (kg)}$.\\
\textbf{Đáp án của câu hỏi là B}
}
\end{ex}

\begin{ex} %[5D7V40]
Một đội công nhân dự kiến hoàn thành công trình trong $60$ ngày. Nhờ tinh thần thi đua, đội đã hoàn thành trước thời hạn $9$ ngày. Đội công nhân đã rút ngắn thời gian so với kế hoạch là bao nhiêu phần trăm?
\choice
{$12\%$}
{\True $15\%$}
{$18\%$}
{$20\%$}
\loigiai{
Tỉ số phần trăm thời gian rút ngắn là:
$9 : 60 = 0{,}15 = 15\%$.\\
\textbf{Đáp án của câu hỏi là B}
}
\end{ex}

\begin{ex} %[5D7V40]
Trong $400\text{ g}$ dung dịch nước muối có chứa $20\text{ g}$ muối tinh. Tỉ số phần trăm của muối trong dung dịch đó là
\choice
{$4\%$}
{\True $5\%$}
{$6\%$}
{$8\%$}
\loigiai{
Tỉ số phần trăm muối trong dung dịch là:
$20 : 400 = 0{,}05 = 5\%$.\\
\textbf{Đáp án của câu hỏi là B}
}
\end{ex}

\begin{ex} %[5D7V40]
Bán một món hàng với giá $120\,000$ đồng thì người bán hàng được lãi $20\,000$ đồng. Tỉ số phần trăm giữa tiền lãi và tiền vốn là
\choice
{$16{,}67\%$}
{\True $20\%$}
{$25\%$}
{$15\%$}
\loigiai{
Tiền vốn của món hàng là:
$120\,000 - 20\,000 = 100\,000\text{ (đồng)}$.\\
Tỉ số phần trăm tiền lãi so với tiền vốn là:
$20\,000 : 100\,000 = 0{,}2 = 20\%$.\\
\textbf{Đáp án của câu hỏi là B}
}
\end{ex}

\begin{ex} %[5D7V40]
Một hộp bóng có $80$ quả bóng gồm ba màu: $28$ quả bóng xanh, $36$ quả bóng đỏ, còn lại là bóng vàng. Tỉ số phần trăm của số bóng vàng so với tổng số bóng trong hộp là
\choice
{$15\%$}
{\True $20\%$}
{$25\%$}
{$30\%$}
\loigiai{
Số quả bóng vàng trong hộp là:
$80 - (28 + 36) = 16\text{ (quả)}$.\\
Tỉ số phần trăm số bóng vàng là:
$16 : 80 = 0{,}2 = 20\%$.\\
\textbf{Đáp án của câu hỏi là B}
}
\end{ex}

\textbf{Phần 2. Trắc nghiệm đúng sai}
\setcounterpageref{ex}{0}

\begin{ex} %[5D7V40]
Theo kế hoạch năm, một phân xưởng dệt phải hoàn thành $1\,200\text{ m}$ vải. Thực tế phân xưởng đã dệt được $1\,380\text{ m}$ vải.
\choiceTF
{\True Phân xưởng đã dệt vượt mức $180\text{ m}$ vải so với kế hoạch}
{\True Phân xưởng đã thực hiện được $115\%$ kế hoạch}
{\True Phân xưởng đã vượt mức kế hoạch là $15\%$}
{Nếu muốn đạt $120\%$ kế hoạch thì phân xưởng phải dệt được tất cả $1\,500\text{ m}$ vải}
\loigiai{
\begin{itemchoice}
\item Đúng. Số mét vải vượt mức: $1\,380 - 1\,200 = 180\text{ (m)}$.
\item Đúng. Tỉ số phần trăm thực hiện: $1\,380 : 1\,200 = 1{,}15 = 115\%$.
\item Đúng. Số phần trăm vượt mức: $115\% - 100\% = 15\%$.
\item Sai. Số mét vải cần để đạt $120\%$ là: $1\,200 \times 120\% = 1\,440\text{ (m)} 
eq 1\,500\text{ m}$.
\end{itemchoice}
}
\end{ex}

\begin{ex} %[5D7V40]
Một mảnh vườn có diện tích $800\text{ m}^2$, trong đó diện tích trồng rau là $320\text{ m}^2$, diện tích trồng hoa là $280\text{ m}^2$, phần diện tích còn lại dùng làm lối đi.
\choiceTF
{\True Diện tích trồng rau chiếm $40\%$ diện tích mảnh vườn}
{\True Diện tích trồng hoa chiếm $35\%$ diện tích mảnh vườn}
{\True Diện tích làm lối đi là $200\text{ m}^2$}
{Diện tích làm lối đi chiếm $20\%$ diện tích mảnh vườn}
\loigiai{
\begin{itemchoice}
\item Đúng. Tỉ số phần trăm diện tích trồng rau: $320 : 800 = 0{,}4 = 40\%$.
\item Đúng. Tỉ số phần trăm diện tích trồng hoa: $280 : 800 = 0{,}35 = 35\%$.
\item Đúng. Diện tích lối đi: $800 - (320 + 280) = 200\text{ (m}^2\text{)}$.
\item Sai. Tỉ số phần trăm lối đi: $200 : 800 = 0{,}25 = 25\% 
eq 20\%$.
\end{itemchoice}
}
\end{ex}

\textbf{Phần 3. Trắc nghiệm trả lời ngắn}
\setcounterpageref{ex}{0}

\begin{ex} %[5D7V40]
Lớp 5B có $32$ học sinh hoàn thành xuất sắc nhiệm vụ học tập trên tổng số $40$ học sinh cả lớp. Tính tỉ số phần trăm học sinh hoàn thành xuất sắc nhiệm vụ của lớp 5B.
\shortans{80}
\loigiai{
Tỉ số phần trăm là:
$32 : 40 = 0{,}8 = 80\%$.
}
\end{ex}

\begin{ex} %[5D7V40]
Một xưởng mộc theo kế hoạch phải đóng $200$ bộ bàn ghế. Thực tế xưởng đã đóng được $240$ bộ bàn ghế. Hỏi xưởng đó đã vượt mức kế hoạch bao nhiêu phần trăm?
\shortans{20}
\loigiai{
Số bộ bàn ghế vượt mức kế hoạch là:
$240 - 200 = 40\text{ (bộ)}$.\\
Số phần trăm vượt mức kế hoạch là:
$40 : 200 = 0{,}2 = 20\%$.
}
\end{ex}

\begin{ex} %[5D7V40]
Một quyển sách có giá bìa là $50\,000$ đồng. Nhân ngày hội sách, quyển sách được bán với giá $40\,000$ đồng. Hỏi quyển sách đã được giảm giá bao nhiêu phần trăm?
\shortans{20}
\loigiai{
Số tiền giảm giá là:
$50\,000 - 40\,000 = 10\,000\text{ (đồng)}$.\\
Số phần trăm giảm giá là:
$10\,000 : 50\,000 = 0{,}2 = 20\%$.
}
\end{ex}

\begin{ex} %[5D7V40]
Một thùng chứa $50\text{ l}$ xăng. Người ta đã lấy ra dùng hết $35\text{ l}$ xăng. Hỏi lượng xăng đã dùng chiếm bao nhiêu phần trăm lượng xăng ban đầu trong thùng?
\shortans{70}
\loigiai{
Tỉ số phần trăm lượng xăng đã dùng là:
$35 : 50 = 0{,}7 = 70\%$.
}
\end{ex}

\begin{ex} %[5D7V40]
Một trang trại nuôi $600$ con vịt và $400$ con gà. Hỏi số con gà chiếm bao nhiêu phần trăm tổng số con gia cầm của trang trại đó?
\shortans{40}
\loigiai{
Tổng số gia cầm là:
$600 + 400 = 1\,000\text{ (con)}$.\\
Số con gà chiếm số phần trăm là:
$400 : 1\,000 = 0{,}4 = 40\%$.
}
\end{ex}

\begin{ex} %[5D7V40]
Một người mua một chiếc xe đạp với giá vốn là $1\,200\,000$ đồng rồi bán lại với giá $1\,500\,000$ đồng. Hỏi số tiền lãi bằng bao nhiêu phần trăm tiền vốn?
\shortans{25}
\loigiai{
Số tiền lãi là:
$1\,500\,000 - 1\,200\,000 = 300\,000\text{ (đồng)}$.\\
Tỉ số phần trăm tiền lãi so với tiền vốn là:
$300\,000 : 1\,200\,000 = 0{,}25 = 25\%$.
}
\end{ex}
